\documentclass[12pt,a4paper]{article}
\usepackage[utf8]{inputenc}
\usepackage[T1]{fontenc}
\usepackage[brazil]{babel}
\usepackage{mathptmx}         
\usepackage{amsmath,amssymb}
\usepackage{booktabs}
\usepackage{array}
\usepackage{float}
\usepackage{graphicx}
\usepackage{setspace}
\usepackage{titlesec}
\usepackage{fancyhdr}
\usepackage[font=small,labelsep=endash]{caption}
\usepackage[a4paper,top=3cm,bottom=2cm,left=3cm,right=2cm]{geometry}
\graphicspath{{./}{figuras/}}
\onehalfspacing
\setlength{\parindent}{1.25cm}
\setlength{\parskip}{0pt}
\pagestyle{fancy}
\fancyhf{}
\fancyhead[R]{\thepage}
\renewcommand{\headrulewidth}{0pt}
\setcounter{secnumdepth}{4}
\setcounter{section}{3}
\renewcommand{\thesection}{\arabic{section}}
\renewcommand{\thesubsection}{\thesection.\arabic{subsection}}
\renewcommand{\thesubsubsection}{\thesubsection.\arabic{subsubsection}}
\renewcommand{\theparagraph}{\thesubsubsection.\arabic{paragraph}}
\titleformat{\section}{\normalfont\normalsize\bfseries\MakeUppercase}{\thesection}{1em}{}
\titleformat{\subsection}{\normalfont\normalsize\bfseries}{\thesubsection}{1em}{}
\titleformat{\subsubsection}{\normalfont\normalsize}{\thesubsubsection}{1em}{}
\titleformat{\paragraph}[block]{\normalfont\normalsize\itshape}{\theparagraph}{1em}{}
\titlespacing*{\section}{0pt}{1.5em}{1em}
\titlespacing*{\subsection}{0pt}{1.5em}{1em}
\titlespacing*{\subsubsection}{0pt}{1.2em}{0.8em}
\titlespacing*{\paragraph}{0pt}{1em}{0.6em}

\renewcommand{\arraystretch}{1.2}

\begin{document}

\section{Análise Estatística do Sistema de Compartilhamento de Bicicletas}

\subsection{Análise da Série Temporal de \texttt{total\_user}}

\subsubsection{Introdução}

O primeiro item da Questão 4 tem como objetivo analisar o comportamento temporal
da utilização do sistema de compartilhamento de bicicletas por meio da variável
\texttt{total\_user}. Para isso, foi construída uma série temporal considerando
as observações diárias do conjunto de dados específico do grupo, formado por 300
registros consecutivos.

A variável \texttt{total\_user} representa o número total de usuários do sistema
em cada dia e é obtida pela soma entre o número de usuários casuais e registrados:

\begin{equation}
\text{total\_user}_i =
\text{casual}_i + \text{registered}_i,
\end{equation}

em que $i$ representa cada observação diária.

A análise da série temporal permite identificar períodos de maior e menor
utilização do sistema, além de observar tendências e oscilações ao longo do
período analisado.

\subsubsection{Construção da série temporal}

Para a construção da série temporal, as observações foram organizadas em ordem
cronológica utilizando a variável \texttt{dteday}. No eixo horizontal do gráfico
foram representadas as datas das observações, enquanto no eixo vertical foi
representado o número total de usuários em cada dia.

A Figura \ref{fig:serie-temporal} apresenta a evolução diária da variável
\texttt{total\_user} durante o período analisado.

\begin{figure}[H]
    \centering
    \includegraphics[width=0.95\textwidth]{4.1.serie_temporal_total_user.png}
        \caption{Série temporal do número total de usuários do sistema.}
    \label{fig:serie-temporal}
\end{figure}

A representação gráfica evidencia variações diárias importantes na utilização
do sistema, além de mudanças no nível geral de demanda ao longo dos meses.

\subsubsection{Identificação dos valores máximo e mínimo}

A partir da série temporal, foram identificadas as observações correspondentes à
maior e à menor utilização do sistema no conjunto analisado.

\begin{table}[H]
\centering
\caption{Maior e menor utilização observadas no período}
\begin{tabular}{lcc}
\toprule
\textbf{Situação} & \textbf{Data} & \textbf{Total de usuários} \\
\midrule
Maior utilização & 04/07/2011 & 6043 \\
Menor utilização & 27/01/2011 & 431 \\
\bottomrule
\end{tabular}
\end{table}

O maior valor registrado ocorreu em 04 de julho de 2011, quando foram
contabilizados 6043 usuários. Em contraste, o menor valor do período foi
observado em 27 de janeiro de 2011, com 431 usuários.

\subsubsection{Análise dos períodos de maior e menor utilização}

\paragraph{Período de menor utilização}

Nos primeiros meses da série, especialmente durante janeiro e fevereiro,
observam-se níveis de utilização relativamente menores. Grande parte das
observações desse período encontra-se em patamares inferiores aos verificados
nos meses posteriores.

O menor valor de toda a série ocorre em 27 de janeiro, com 431 usuários. Além
desse ponto mínimo, o início da série apresenta diversos dias com valores
próximos ou inferiores a aproximadamente 2000 usuários, indicando um período de
demanda reduzida em comparação com os meses seguintes.

\paragraph{Período de crescimento da utilização}

Entre março e abril, a série apresenta uma elevação gradual no número total de
usuários. Os valores passam progressivamente de níveis próximos de 2000 usuários
para observações frequentemente superiores a 3000 e 4000 usuários.

Esse comportamento caracteriza uma mudança importante no nível de utilização do
sistema. Embora permaneçam algumas quedas pontuais, o comportamento geral indica
aumento da demanda ao longo desse intervalo.

\paragraph{Período de maior utilização}

A partir de maio, a série passa a apresentar valores predominantemente mais
elevados. Entre maio e julho, grande parte das observações encontra-se entre
aproximadamente 4000 e 5500 usuários.

O ponto máximo da série ocorre em 04 de julho de 2011, com 6043 usuários. Esse
resultado está inserido em um período no qual vários dias também apresentam
valores próximos ou superiores a 5000 usuários, caracterizando uma fase de
utilização elevada do sistema.

\paragraph{Comportamento após o valor máximo}

Após julho, o número total de usuários permanece elevado em muitos dias, porém
a série passa a apresentar oscilações mais intensas.

Entre agosto e outubro, são observados diversos valores entre aproximadamente
4000 e 5000 usuários, intercalados com quedas acentuadas. Dessa forma, embora o
nível geral de utilização permaneça elevado, a variabilidade diária torna-se
mais evidente.

\subsubsection{Principais padrões observados}

A análise da série temporal permite identificar três características principais.
Primeiramente, observa-se uma tendência de crescimento da utilização durante a
primeira metade do período analisado. Os valores registrados em janeiro e
fevereiro são, em geral, inferiores aos observados nos meses seguintes, enquanto
março e abril representam uma fase de transição para níveis mais elevados de
demanda.

Em segundo lugar, entre maio e julho ocorre o período de utilização mais intensa.
Nesse intervalo, os valores permanecem frequentemente acima de 4000 usuários,
culminando no máximo de 6043 usuários em 04 de julho.

Por fim, após o período de crescimento, a série apresenta um nível elevado de
utilização acompanhado de maior variabilidade. Entre agosto e outubro são
observadas sucessivas oscilações, com dias de utilização elevada intercalados
com quedas acentuadas.

\subsubsection{Conclusão}

A construção da série temporal de \texttt{total\_user} permitiu visualizar a
evolução da utilização do sistema de compartilhamento de bicicletas ao longo das
300 observações do grupo.

O menor valor observado foi de 431 usuários em 27 de janeiro de 2011, enquanto
o maior valor foi de 6043 usuários em 04 de julho de 2011. A série apresenta
níveis relativamente baixos no início do período, seguidos por crescimento
progressivo durante março e abril.

Entre maio e julho, a utilização atinge seus maiores níveis, com grande
concentração de observações acima de 4000 usuários. Após esse período, a demanda
permanece relativamente elevada, porém acompanhada por oscilações diárias mais
acentuadas.

Dessa forma, a análise temporal evidencia tanto uma tendência de crescimento na
primeira metade do período quanto uma maior variabilidade na segunda metade,
fornecendo uma visão geral da evolução da demanda pelo sistema ao longo do tempo.

\subsection{Associação entre Temperatura e Utilização do Sistema}

Para aprofundar a análise iniciada na Questão 3 a respeito do impacto térmico, selecionou-se a variável quantitativa contínua da temperatura (\texttt{temp}) para avaliar estatisticamente a sua associação com o volume diário de usuários (\texttt{total\_user}).

Como ambas as variáveis são contínuas, a medida estatística adequada escolhida foi o \textbf{Coeficiente de Correlação de Pearson}. Esta métrica quantifica a força e a direção da relação linear entre os dados. Para a representação visual, optou-se por um gráfico de dispersão (\textit{scatter plot}) acrescido de uma reta de regressão linear para facilitar a observação da tendência.

A execução do teste estatístico computacional resultou em um coeficiente de correlação de \textbf{0,8058}. O comportamento visual desta distribuição é apresentado na Figura \ref{fig:dispersao_temp}.

\begin{figure}[H]
    \centering
    \caption{Associação linear entre temperatura e total de usuários diários.}
    \vspace{0.1cm}
    \includegraphics[width=0.8\textwidth]{4.2.grafico_dispersao_temp.png}
    \\ \vspace{0.1cm}
    \footnotesize{Fonte: Elaborado pelos autores (2026) com base nos dados de \texttt{2.1.total\_user.xlsx}.}
    \label{fig:dispersao_temp}
\end{figure}

\textbf{Interpretação e Comparação com Resultados Anteriores}

O valor obtido ($r \approx 0,81$) indica uma correlação linear positiva e de forte intensidade. Isso significa que há uma proporção direta bastante consistente: a elevação da temperatura está intimamente associada ao crescimento do número total de ciclistas no sistema.

Este resultado fundamenta matematicamente as conclusões qualitativas obtidas na Questão 3. Naquela etapa, constatou-se que os dias com temperaturas mais agradáveis e quentes (como os meses centrais do ano, cujo ápice de uso ocorreu em julho) concentravam maiores volumes de demanda. O teste de Pearson comprova de forma contundente que esse comportamento não é apenas uma flutuação sazonal, mas sim uma tendência estatística robusta. A temperatura atua como um forte preditor do volume de aluguéis, validando plenamente as hipóteses levantadas na etapa anterior do estudo.

\subsection{Associação entre Condições Meteorológicas e Utilização do Sistema}

Para a segunda característica selecionada na Questão 3, as condições meteorológicas (\texttt{weathersit}), a abordagem estatística requerida é diferente. Por se tratar de uma variável categórica que divide os dias em grupos distintos (1 = Bom, 2 = Nublado/Misto, 3 = Chuva/Neve Leve) a ser comparada com uma variável contínua (\texttt{total\_user}), utilizou-se a \textbf{Análise de Variância (ANOVA)}. Para a visualização da distribuição dos dados em cada categoria, optou-se pelo \textit{boxplot}.

A execução do teste ANOVA devolveu um \textit{p-value} de aproximadamente $2,13 \times 10^{-9}$. A representação visual desta distribuição encontra-se na Figura \ref{fig:boxplot_clima}.

\begin{figure}[H]
    \centering
    \caption{Associação entre condições meteorológicas e total de utilizadores diários.}
    \vspace{0.1cm}
    \includegraphics[width=0.8\textwidth]{4.3.boxplot_clima.png}
    \\ \vspace{0.1cm}
    \footnotesize{Fonte: Elaborado pelos autores (2026) com base nos dados de \texttt{2.1.total\_user.xlsx}.}
    \label{fig:boxplot_clima}
\end{figure}

\textbf{Interpretação e Comparação com Resultados Anteriores}

O \textit{p-value} obtido ($p < 0,05$) indica que existe uma diferença estatisticamente significativa na média de utilizadores entre os diferentes tipos de clima. Observando o \textit{boxplot}, é evidente que a mediana e a concentração de alugueres são consideravelmente superiores nos dias de clima limpo, sofrendo uma queda drástica nos dias de precipitação.

Esta validação matemática corrobora de forma absoluta as observações qualitativas documentadas na Questão 3. A redução do uso do sistema não é um evento aleatório, mas sim uma consequência direta e estatisticamente provada do agravamento das condições meteorológicas, consolidando o clima como um fator determinante na procura por bicicletas.

\subsection{Temperatura e baixa utilização}

\subsubsection{Critério e preparação dos dados}

A análise mantém a seleção definida na Questão 1: com $M=581706$,
$r=1+(M\bmod 100)=7$, sendo utilizadas as $300$ observações consecutivas
de 07/01/2011 a 02/11/2011. O total diário e a classificação seguem as
definições da Questão 2:
\[
U_i=\texttt{casual}_i+\texttt{registered}_i,
\qquad Q_1(U)=2105{,}5,
\]
\[
\texttt{low\_usage}_i=
\begin{cases}
1, & U_i<Q_1(U),\\
0, & U_i\geq Q_1(U).
\end{cases}
\]
O quartil foi calculado sobre os $300$ dias pelo método
\texttt{quantile(type = 7)} do R, antes de dividir as observações em grupos.
Assim, $75$ dias ($25\%$) são classificados como de baixa utilização e os
$225$ restantes pertencem ao outro grupo. Não se recalcula um quartil para
cada grupo ou estação.

Para tornar o eixo horizontal interpretável, a temperatura normalizada do
enunciado é convertida por $T_{\mathrm{C}}=41T_{\mathrm{norm}}$. O script
também reconhece uma coluna que já esteja em graus Celsius. Nenhuma
observação é excluída ou reordenada.

\subsubsection{Dispersão e classificação dos dias}

A Figura~\ref{fig:dispersao} representa a temperatura e o total de usuários
em cada dia. A linha horizontal marca o mesmo $Q_1$ usado na classificação:
os pontos abaixo dela pertencem a \texttt{low\_usage = 1}. A concentração
dos dias de baixa utilização ocorre predominantemente em temperaturas
menores, enquanto os totais mais elevados são observados com maior
frequência em temperaturas moderadas a altas.

\begin{figure}[H]
\centering
\includegraphics[width=0.84\linewidth]{questao_4_3_dispersao.png}
\caption{Temperatura e total diário, com os $300$ dias diferenciados pela
classificação de baixa utilização. A linha tracejada é $Q_1=2105{,}5$.
Gráfico produzido pelo script \texttt{questao\_4\_3.R}.}
\label{fig:dispersao}
\end{figure}

\subsubsection{Comparação visual entre os grupos}

A Figura~\ref{fig:paineis} separa os dois grupos, preservando as mesmas
escalas nos eixos. As retas descrevem a tendência linear \emph{dentro de
cada subconjunto} e foram desenhadas apenas no intervalo de temperatura
observado no respectivo grupo. A comparação mostra uma tendência positiva
aparentemente menos marcada entre os dias de baixa utilização.

\begin{figure}[H]
\centering
\includegraphics[width=0.80\linewidth]{questao_4_3_paineis.png}
\caption{Temperatura e utilização nos dois grupos. As escalas são iguais;
as retas representam ajustes lineares descritivos. Gráfico produzido pelo
script \texttt{questao\_4\_3.R} a partir da tabela do PDF, cuja temperatura
é arredondada. Execute o script com o CSV original para os valores finais.}
\label{fig:paineis}
\end{figure}

\begin{table}[H]
\centering
\small
\caption{Temperatura e correlação de Pearson em cada grupo.}
\begin{tabular}{@{}lrrrr@{}}
\toprule
Grupo & Dias & Média de $T$ & Intervalo de $T$ & $r(T,U)$ \\
\midrule
Baixa utilização & 75 & $11{,}08\,^{\circ}\mathrm C$ & $2{,}4$--$27{,}9\,^{\circ}\mathrm C$ & $0{,}277$ \\
Demais dias & 225 & $24{,}50\,^{\circ}\mathrm C$ & $10{,}8$--$34{,}8\,^{\circ}\mathrm C$ & $0{,}568$ \\
\bottomrule
\end{tabular}
\end{table}

\subsubsection{Interpretação e limitações}

O coeficiente de Pearson é positivo em ambos os grupos, mas menor entre os
dias classificados como de baixa utilização ($r\approx0{,}277$) do que nos
demais ($r\approx0{,}568$). As temperaturas médias também diferem:
$11{,}08\,^{\circ}\mathrm C$ e $24{,}50\,^{\circ}\mathrm C$,
respectivamente. A presença de dias quentes com pouco uso impede uma leitura
determinista da relação: em 27/08/2011, por exemplo, constam aproximadamente
$27{,}9\,^{\circ}\mathrm C$ e $1115$ usuários.

Essa comparação é descritiva. A variável \texttt{low\_usage} foi construída
a partir do próprio total de usuários, representado no eixo vertical. A
separação restringe por definição a faixa de $U$ em cada grupo; diferenças
entre correlações ou inclinações calculadas depois do corte não demonstram
efeitos causais distintos da temperatura. Estação do ano e condições
meteorológicas também podem estar associadas à utilização. Na amostra,
$60$ dos $75$ dias de baixa utilização ocorreram no inverno.

Nas dez primeiras observações do grupo, todos os dias satisfazem
$U_i<2105{,}5$. Portanto, não existe nesse recorte inicial uma segunda
classe com a qual comparar a relação. A avaliação entre os dois grupos
depende das $300$ observações completas.

\subsection{Conclusão}

A análise das 300 observações revela que a utilização do sistema de bicicletas é severamente influenciada por fatores sazonais e meteorológicos. A série temporal (Questão 4) ilustra bem a evolução da procura, registando um mínimo no início do ano (431 utilizadores) e atingindo o pico absoluto nos meses centrais (6.043 utilizadores).

Dentre as características avaliadas na Questão 3, a temperatura e a estação do ano destacaram-se. Observou-se uma forte correlação positiva entre a temperatura e a utilização do sistema ($r \approx 0,81$). O inverno concentrou os períodos de maior ociosidade, com 82,19\% dos seus dias situados abaixo do primeiro quartil global ($Q_1 = 2105,5$, Questão 2).

A Questão 4 demonstrou que a restrição de alugueres diários enfraquece a correlação com a temperatura, decrescendo de $r \approx 0,568$ (uso normal) para $r \approx 0,277$ (baixa utilização). Cruzando os dados climáticos, constatou-se que a precipitação reduz a média diária de alugueres de 3.851 (céu limpo) para apenas 1.562. Estes resultados indicam que a gestão operacional da frota deve alinhar-se rigorosamente às previsões meteorológicas diárias.

\end{document}
